% =============================================================================
%  pretextual/resumo.tex
%  Contém: Resumo (PT-BR) e Abstract (EN)
%
%  ATENÇÃO: O resumo é o ÚLTIMO elemento a ser escrito no TCC.
%  Só escreva quando o trabalho estiver finalizado.
%
%  ESTRUTURA DO PARÁGRAFO ÚNICO:
%    1. Contexto / Problema → Por que o trabalho existe?
%    2. Solução / Método   → O que foi desenvolvido e como?
%    3. Resultados         → O que foi alcançado?
%
%  Tamanho recomendado: 150 a 500 palavras (verifique a norma da sua escola).
% =============================================================================


% -----------------------------------------------------------------------------
%  RESUMO — em português
% -----------------------------------------------------------------------------
\begin{resumo}
  TODO: [Contexto/Problema] Descreva o cenário que motivou o trabalho e
  qual lacuna ou dificuldade ele resolve. [Solução] Explique o que foi
  desenvolvido, citando as principais tecnologias e a abordagem adotada.
  [Resultado] Apresente os resultados obtidos e o impacto ou valor gerado
  pela solução. Lembre-se: um único parágrafo, sem tópicos ou subdivisões.

  % \vspace{\onelineskip} adiciona uma linha em branco antes das palavras-chave
  \vspace{\onelineskip}
  \noindent
  \textbf{Palavras-chave}: TODO: palavra-chave 1; palavra-chave 2;
  palavra-chave 3; palavra-chave 4;
  % Separadas por ponto, em letras minúsculas, de 3 a 6 termos.
\end{resumo}


% -----------------------------------------------------------------------------
%  ABSTRACT — em inglês
%  Tradução fiel do resumo acima.
%  O ambiente otherlanguage* muda a hifenização para inglês só neste trecho.
% -----------------------------------------------------------------------------
\begin{resumo}[Abstract]
  \begin{otherlanguage*}{english}
    TODO: [Context/Problem] Describe the scenario that motivated this work
    and the gap or difficulty it addresses. [Solution] Explain what was
    developed, mentioning the main technologies and the approach adopted.
    [Results] Present the results achieved and the impact or value generated
    by the solution.

    \vspace{\onelineskip}
    \noindent
    \textbf{Keywords}: TODO: keyword 1; keyword 2; keyword 3; keyword 4;
  \end{otherlanguage*}
\end{resumo}
